\section{Validation Against Observed Process Execution}
\label{app:eventlog_validation}

Appendix~\ref{app:apqc_validation} scores our ordering instrument against sequences that practitioners
\emph{documented}, and concedes a limitation it cannot overcome: a documented ordering is an editorial
convention, and for branches whose children are parallel or mutually exclusive no true sequence exists, so
disagreement there says nothing about whether the instrument understands sequence.
This appendix removes that limitation by moving from documented order to \emph{observed} order.

A process-mining event log records, case by case, the activities an organization actually executed and when.
Two things follow.
First, the ordering is a fact about production rather than an assertion about it.
Second, and more useful, the data reveal \emph{which precedences are real}: if activity $x$ precedes $y$ in
97\% of cases the precedence is genuine, whereas a 50/50 split means the two run concurrently and no
ordering could be correct.
We can therefore condition on the pairs that have a true order before asking whether the model recovers it,
which is precisely the step the documented benchmark does not permit.

\subsection{Data}
\label{app:eventlog_data}

We use eight publicly available event logs from the 4TU.ResearchData repository, listed in
Table~\ref{tab:eventlog_overview}: fine collection by a local police force, sepsis care in a hospital
emergency department, consumer lending at a financial institution, and five university travel and
reimbursement processes.
Together they cover 215,715 cases and 2,048,949 recorded events.
Their activity sets are considerably larger than the APQC process branches, between 8 and 28 activities
against a median of five, so they also test the instrument on long sequences.

From each log we recover one trace per case, the sequence of activity labels in timestamp order, and keep the
activities appearing in at least 2\% of cases so that the target is not driven by rare exceptions.
We use each activity's \emph{first} occurrence within a case, which handles the loops present in every real
log such as repeated laboratory tests or repeated approval rounds.
This yields the two objects we need.
The \emph{observed ordering} ranks activities by their mean normalized first-occurrence position across
cases; this is the analogue of APQC's published order.
The \emph{pairwise precedence matrix} records, for every pair, the share of cases containing both in which
one precedes the other.
We call a pair \emph{determinate} when one direction holds in at least 80\% of cases.
Across the eight logs, 1,124 of 1,178 activity pairs (95.4\%) are determinate, so the benchmark is a
statement about sequence for almost all of the comparisons we make.

\begin{table}[h!]
    \caption{Event Logs Used for Validation}
    \label{tab:eventlog_overview}
    \begin{center}
    \small
    \begin{tabular}{llrrrr}
    \toprule
    Log & Process & Cases & Events & Activities & Determinate pairs \\
    \midrule
    Road traffic fines & Fine management, local police & 150,370 & 561,470 & 8 & 27 of 28 \\
    Sepsis & Emergency department sepsis care & 1,050 & 15,214 & 15 & 91 of 98 \\
    BPI 2017 & Loan application handling & 31,509 & 1,202,267 & 23 & 238 of 245 \\
    BPI 2020 & Domestic travel declarations & 10,366 & 56,303 & 10 & 41 of 44 \\
    BPI 2020 & International travel declarations & 6,449 & 72,151 & 22 & 218 of 224 \\
    BPI 2020 & Requests for payment & 6,814 & 36,724 & 10 & 39 of 42 \\
    BPI 2020 & Prepaid travel costs & 2,092 & 18,239 & 17 & 128 of 129 \\
    BPI 2020 & Travel permits and declarations & 7,065 & 86,581 & 28 & 342 of 368 \\
    \midrule
    Total & & 215,715 & 2,048,949 & 133 & 1,124 of 1,178 \\
    \bottomrule
    \end{tabular}
    \end{center}
    \raggedright \footnotesize{\emph{Notes:}} All logs are openly available from the 4TU.ResearchData repository.
    ``Activities'' counts those appearing in at least 2\% of a log's cases, which is the set handed to the model.
    A pair of activities is ``determinate'' when one of the two orders holds in at least 80\% of the cases containing both.
\end{table}

\subsection{Design}
\label{app:eventlog_design}

The procedure mirrors Appendix~\ref{app:apqc_validation_design} exactly.
For each log we shuffle the activity labels under a fixed seed, so the model never sees the observed
sequence, and ask it to order them, supplying a one-line description of the process in place of an
occupation title.
We use all eleven prompts, the main specification of Appendix~\ref{app:prompts} and the ten alternatives of
Appendix~\ref{app:gptPrompts_robustness}, carried over with the same mechanical substitution that maps the
occupation domain onto the process domain.
No log is told which activity starts or ends a case.

We report two levels of analysis.
At the \emph{log level} we compute Kendall's $\tau$ between the model's ordering and the observed ordering,
which is directly comparable to the APQC numbers, benchmarked against a permutation null of $1{,}000$ random
orderings per log.
At the \emph{pair level} we ask only whether the model places a pair in the direction the traces show.
The pair level is the more informative of the two here, both because it is where the determinacy
conditioning bites and because it is closer to what the paper's predictions actually use: Predictions~\#1
and~\#3 depend on adjacency rather than on the global ordering.

\subsection{Results}
\label{app:eventlog_results}

Table~\ref{tab:eventlog_ordering_validation} reports the results and
Figure~\ref{fig:eventlog_activity_accuracy} the underlying distribution.

Against a null of $0.5$, the model places $78.8\%$ of all activity pairs in the direction the traces show.
Restricting to the determinate pairs, those the data confirm have a real order, gives $78.9\%$, and weighting
each pair by how determinate it is gives $78.8\%$.
That these three numbers coincide is itself informative: the model does no better on the pairs that a real
organization orders strictly than on pairs in general, so the aggregate is not being propped up by easy
cases.
At the log level the mean Kendall's $\tau$ is $0.48$ with a median of $0.56$, against a permutation null
centered at $0.00$ ($z = 6.8$, $p = 0.001$).
The model is not echoing the list it was shown: the average $\tau$ between its output and the shuffled input
is $0.007$.

The ten alternative prompts again leave the picture unchanged.
Determinate-pair accuracy under individual prompts spans $0.767$ to $0.817$ and log-level mean $\tau$ spans
$0.433$ to $0.559$, so no conclusion here depends on the phrasing used in the main analysis.

The result that matters for the paper is the comparison with Appendix~\ref{app:apqc_validation}.
Restricting both exercises to \emph{adjacent} pairs, the statistic they share, the documented benchmark gives
$0.700$ under the main prompt and $0.697$ averaged over the alternatives, while the observed benchmark gives
$0.679$ and $0.699$.
Two independent benchmarks, one written down by practitioners and one recovered from what organizations
actually did, agree to within roughly two percentage points about how much sequence structure the instrument
recovers.

\begin{table}[h!]
    \caption{Overlap Between GPT-Generated Orderings and Observed Process Execution}
    \label{tab:eventlog_ordering_validation}
    \begin{center}
    \input{tables/eventlog_ordering_validation.tex}
    \end{center}
    \raggedright \footnotesize{\emph{Notes:}} Pairwise precedence accuracy is the share of activity pairs the model places in the direction the event traces show, with a null value of $0.5$.
    ``Determinate pairs'' are those for which one order holds in at least 80\% of the cases containing both activities.
    ``Adjacent pairs'' are consecutive in the observed ordering and are the statistic directly comparable to the adjacent-pair agreement of Table~\ref{tab:apqc_pcf_ordering_validation}.
    Column~(2) pools the ten alternative prompts of Appendix~\ref{app:gptPrompts_robustness}, so its pair counts are ten times those of column~(1).
    The random-ordering row is the mean of $1{,}000$ random permutations per log.
\end{table}

\begin{figure}[h!]
\caption{Distribution of Overlap With Observed Execution Across Activities}
\label{fig:eventlog_activity_accuracy}
\begin{center}
\begin{subfigure}[b]{0.49\textwidth}
    \includegraphics[width=\linewidth]{plots/eventlog_ordering/activity_accuracy_main_share.png}
    \caption{Main prompt}
    \label{fig:eventlog_acc_main}
\end{subfigure}
\hfill
\begin{subfigure}[b]{0.49\textwidth}
    \includegraphics[width=\linewidth]{plots/eventlog_ordering/activity_accuracy_alt_share.png}
    \caption{Average of ten alternative prompts}
    \label{fig:eventlog_acc_alt}
\end{subfigure}
\end{center}
\footnotesize{
\emph{Notes:} The unit is an activity (133 across the eight logs). For each activity we take the determinate pairs it belongs to and plot the share the model ordered as the traces show.
The black dashed line marks the coin flip and the red dashed line the mean.
Averaging over prompts raises the share of activities above the coin flip from $86.5\%$ to $91.7\%$ while leaving the mean almost unchanged, the same pattern as in Figure~\ref{fig:apqc_tau_distribution}.
Because activities within a log share one ordering and each pair contributes to two activities, this distribution is descriptive; the inferential statistics are the pooled pair-level accuracy and the log-level permutation test.
}
\end{figure}

\subsection{Where the Instrument Fails, and Why It Understates Its Own Accuracy}
\label{app:eventlog_failures}

Performance varies sharply across logs, and the variation is diagnostic.
Determinate-pair accuracy is $0.96$ for road traffic fines, $0.95$ for sepsis care and $0.91$ for
international declarations, but $0.56$ for requests for payment and $0.49$ for domestic declarations.
Inspecting the errors, the failures are overwhelmingly about \emph{reading the labels}, not about
understanding sequence.

Three examples carry most of the damage.
In the domestic declaration log the model ranks \texttt{Request Payment} first, reading it as the employee's
initiating request; in the data it is the second-to-last step, the finance system requesting disbursement.
In the loan log the model treats \texttt{A\_Accepted} as terminal, whereas in the data it is an early state
meaning the application entered processing, preceding validation, offer creation, and completion; seven of
the twelve most determinate errors are this single activity.
In the sepsis log the model places \texttt{Return ER} mid-visit, while in the data it denotes a patient
returning to the emergency department after discharge and is genuinely last.

This matters for how the number should be read.
Activity labels in an event log are the strings an information system emits, not descriptions of work.
Some are terse (\texttt{A\_Concept}), and many encode the actor rather than the action
(\texttt{Declaration APPROVED by ADMINISTRATION}).
The model is therefore working with substantially less context than an O*NET task statement provides, and the
errors we observe are of a kind that richer task descriptions would largely prevent.
Accuracy here is, in that specific sense, a conservative estimate of what the procedure achieves on the task
statements the paper actually uses.

\subsection{Discussion}
\label{app:eventlog_discussion}

Taken with Appendix~\ref{app:apqc_validation}, this exercise closes the objection that motivated both.
If the orderings were an artifact of language-model reasoning, with semantically similar items placed
adjacent regardless of production logic, there would be no reason for them to reproduce the precedence
structure of 215,715 executed cases in four unrelated domains, and no reason for a documented benchmark and
an observed one to agree so closely on how much structure is recovered.

Two limitations should be stated plainly.
First, the eight logs are not a random sample of work.
They are administrative and clinical processes with recorded case identifiers, which is what makes them
loggable at all, and they under-represent the judgment-intensive occupations where AI exposure is
concentrated.
Second, and as with the APQC exercise, this validates the ordering \emph{procedure} rather than the specific
O*NET orderings used in the paper, since no crosswalk between event-log activities and O*NET tasks is
involved.
What the two appendices jointly establish is that the procedure recovers a substantial and highly
non-random share of real sequence structure, under every prompt formulation we try, against both a
documented and an observed benchmark.
